\chapter{Entrelaceur avec les tenseurs symétriques sans trace et réduction transverse \texorpdfstring{\(12\to5\to5\to2\)}{12→5→5→2}}
\label{chap:parte-iv-64}

Le repère icosaédrique construit une isométrie équivariante vers l’espace des tenseurs symétriques sans trace. La projection transverse réalise \(12\to5\to5\to2\), conserve le support à cinq composantes et localise le sous-espace radiatif de dimension \(2\).

\section{Entrelaceur avec l’espace des tenseurs symétriques sans trace}

Soit \(\mathcal V\subset S^2\) l’ensemble des douze sommets unitaires de l’icosaèdre. Pour chaque \(v\in\mathcal V\), définissons
\begin{equation}
\label{rm:eq:gravity-Qv}
Q_v=vv^T-\frac13I_3\in
\operatorname{STF}_3:=\operatorname{Sym}^2_0(\R^3).
\end{equation}
L’application de repère
\begin{equation}
\label{rm:eq:gravity-Gamma0}
\Gamma_0\!\left(\sum_{v\in\mathcal V}x_ve_v\right)
=\sum_{v\in\mathcal V}x_vQ_v
\end{equation}
est \(A_5\)-équivariante et satisfait
\begin{equation}
\label{rm:eq:gravity-tight-frame}
\Gamma_0\Gamma_0^*=\frac85I_{\operatorname{STF}_3},
\qquad
\Gamma_0=\Gamma_0P_5.
\end{equation}

\begin{theorem}[Entrelaceur tensoriel normalisé]
\label{rm:thm:gravity-gamma5}
L’application
\begin{equation}
\label{rm:eq:gravity-gamma5}
\boxed{
\Gamma_5:=\sqrt{\frac58}\,\Gamma_0|_{V_5}:
V_5\xrightarrow{\ \simeq\ }\operatorname{STF}_3}
\end{equation}
est une isométrie \(A_5\)-équivariante. Son adjoint est
\begin{equation}
\label{rm:eq:gravity-gamma5-adjoint}
\Gamma_5^*Q=\sqrt{\frac58}
\sum_{v\in\mathcal V}(v^TQv)e_v.
\end{equation}
\end{theorem}

\begin{proof}
L’équivariance découle de \(Q_{gv}=gQ_vg^T\). L’opérateur
\(\Gamma_0\Gamma_0^*\) commute avec la représentation irréductible de
\(A_5\) sur \(\operatorname{STF}_3\) ; il est donc scalaire. Sa trace
est
\[
\sum_{v\in\mathcal V}\|Q_v\|_F^2
=12\cdot\frac23=8;
\]
en divisant par la dimension cinq, on obtient \(8/5\). D’où
\(\Gamma_5\Gamma_5^*=I\). Comme le domaine et le codomaine sont de dimension
cinq, on a aussi \(\Gamma_5^*\Gamma_5=I\). La formule de l’adjoint
découle de \(\langle Q,Q_v\rangle_F=v^TQv\).
\end{proof}

La normalisation au moyen du repère icosaédrique élimine une coïncidence purement cardinale : elle construit l’application qui transporte le facteur fini vers le support continu de spin deux.

\section{Projection transversale sans trace}

Pour \(n\in S^2\), soit \(\Pi(n)=I-nn^T\). Nous définissons
\begin{equation}
\label{rm:eq:gravity-lambda-tt}
\boxed{
\Lambda(n)Q
=\Pi Q\Pi-\frac12\Tr(\Pi Q\Pi)\Pi.}
\end{equation}

\begin{theorem}[Réduction de cinq composantes à deux]
\label{rm:thm:gravity-tt}
La restriction de \(\Lambda(n)\) à \(\operatorname{STF}_3\) est le
projecteur orthogonal sur
\begin{equation}
\label{rm:eq:gravity-htt}
\mathcal H_{\rm TT}(n)
=\{Q\in\operatorname{STF}_3:Qn=0\}.
\end{equation}
En particulier,
\begin{equation}
\label{rm:eq:gravity-lambda-properties}
\Lambda(n)^2=\Lambda(n)=\Lambda(n)^*,
\qquad \rank\Lambda(n)=2.
\end{equation}
Pour \(n=e_3\), l’image est engendrée par
\begin{equation}
\label{rm:eq:gravity-plus-cross}
E_+=\frac1{\sqrt2}\operatorname{diag}(1,-1,0),
\qquad
E_\times=\frac1{\sqrt2}
\begin{pmatrix}0&1&0\\1&0&0\\0&0&0\end{pmatrix}.
\end{equation}
\end{theorem}

\begin{proof}
Comme \(\Pi n=0\), la sortie est transversale. Sa trace vaut
\[
\Tr(\Pi Q\Pi)-\frac12\Tr(\Pi)\Tr(\Pi Q\Pi)=0,
\]
car \(\Tr\Pi=2\). Si \(Qn=0\) et \(\Tr Q=0\), alors
\(\Pi Q\Pi=Q\) et \(\Tr(\Pi Q\Pi)=0\) ; donc
\(\Lambda(n)Q=Q\). La formule est autoadjointe pour le produit de
Frobenius. Dans le repère \(n=e_3\), sa matrice dans une base STF adaptée est
\(\operatorname{diag}(1,1,0,0,0)\), ce qui prouve le rang.
\end{proof}

La chaîne complète devient donc
\begin{equation}
\label{rm:eq:gravity-12-5-5-2}
\boxed{
V_{12}\xrightarrow{P_5}V_5
\xrightarrow{\Gamma_5}\operatorname{STF}_3
\xrightarrow{\Lambda(n)}\mathcal H_{\rm TT}(n),
\qquad 12\longrightarrow5\longrightarrow5\longrightarrow2.}
\end{equation}
Le couplage réduit
\begin{equation}
\label{rm:eq:gravity-effective-map}
\mathbb G_{{\rm TT},a}(n)
=G_a(\theta_{108}^{\rm circ})\Lambda(n)\Gamma_5P_5
\end{equation}
est de rang deux.


\paragraph{Décomposition en hélicités du spectre et séparation des modes massif et sans masse.} La réduction permet de distinguer le spectre complet des hélicités d’une réalisation massive et ses deux hélicités sans masse.
